\documentclass[11pt]{article}
\usepackage[margin=1in]{geometry}
\usepackage{amsmath,amssymb,amsthm,mathtools}
\usepackage{graphicx}
\usepackage{hyperref}
\usepackage{xcolor}

\newtheorem{theorem}{Theorem}
\newtheorem{lemma}{Lemma}
\newtheorem{remark}{Remark}
\newcommand{\R}{\mathbb R}
\newcommand{\cl}{\operatorname{cl}}
\newcommand{\dist}{\operatorname{dist}}

\title{A little-Zygmund graph with no normal:\\
a counterexample to Conjecture 6.1 of Balashov--Repov\v{s}}
\author{Autonomous solution packet for arXiv:1007.0162}
\date{August 11, 2026}

\begin{document}
\maketitle

\begin{center}
\fbox{\parbox{0.91\textwidth}{
\textbf{Status: candidate full counterexample, likely valid.}
The construction is elementary and self-contained.  Its only delicate point is a
uniform second-difference estimate, proved below.  No claim of priority is made:
targeted cheap-index and web searches found no resolution of the conjecture, but
the novelty search was not exhaustive.
}}
\end{center}

\section{The question}

For a closed subset $C$ of a normed space, write
\[
 \gamma_C(\varepsilon)=
 \inf\left\{r>0:
 B_r\!\left(\frac{x+y}{2}\right)\cap C\ne\varnothing
 \text{ whenever }x,y\in C, \|x-y\|\leq\varepsilon\right\}.
\]
Thus $C$ is weakly convex on a small scale if
$\gamma_C(\varepsilon)<\varepsilon/2$ there.  Conjecture 6.1 of
Balashov--Repov\v{s} asks, in particular, whether the assumptions
\[
 \gamma(\varepsilon)/\varepsilon\longrightarrow0,qquad
 A=\cl(\operatorname{int}A),
\]
for both $A$ and $\cl(E\setminus A)$ force the existence of a unit normal at
every boundary point, uniformly continuously depending on that point.  The
ambient space is allowed to be uniformly convex and, parenthetically, smooth.

The answer is negative already in the Euclidean plane.

\begin{theorem}\label{thm:main}
There is a globally Lipschitz function $f:\R\to\R$ such that, for
\[
 A=\{(x,y):y\geq f(x)\},\qquad
 B=\cl(\R^2\setminus A)=\{(x,y):y\leq f(x)\},
\]
the following hold.
\begin{enumerate}
 \item $A$ and $B$ are closed and weakly convex, with
 $\gamma_A(\varepsilon)/\varepsilon\to0$ and
 $\gamma_B(\varepsilon)/\varepsilon\to0$.
 \item $A=\cl(\operatorname{int}A)$.
 \item The Fr\'echet normal cone of $A$ at the origin is trivial:
 $N_A(0,0)=\{0\}$.
\end{enumerate}
Consequently $A$ has no unit normal, and no unit proximal normal, at the
origin.  Hence Conjecture 6.1 is false even when $E=\R^2$ is both uniformly
convex and uniformly smooth.
\end{theorem}

\section{Construction and the second-difference estimate}

Fix $0<r_0<1/4$ and an even function
$\chi\in C_c^\infty(\R)$ satisfying
\[
 \chi(x)=1\quad (|x|\leq r_0),\qquad
 \chi(x)=0\quad (|x|\geq2r_0).
\]
Set
\[
 f(0)=0,\qquad
 f(x)=\chi(x)x\sin\!\left(\log\log\frac e{|x|}\right)
 \quad(x\ne0).
\]
The expression is well-defined on the support of $\chi$.  It is globally
Lipschitz: near zero its derivative is bounded, and away from zero it is
smooth and compactly supported.

For $h\in\R$, let
\[
 \Delta_h^2f(x)=f(x+h)+f(x-h)-2f(x),\qquad
 \Omega(t)=\sup_{x\in\R,\ |h|\leq t}|\Delta_h^2f(x)|.
\]

\begin{lemma}[little-Zygmund estimate]\label{lem:zyg}
One has $\Omega(t)=o(t)$ as $t\downarrow0$.
\end{lemma}

\begin{proof}
It suffices to take $h>0$.  Put
\[
 a(r)=\sin\!\left(\log\log\frac e r\right).
\]
In the core $0<|x|\leq r_0$, direct differentiation gives
\[
 |f''(x)|\leq
 \frac{C}{|x|\log(e/|x|)}.                                      \tag{1}
\]
If $|x|\geq2h$ and $x-h,x,x+h$ remain in the core, symmetric Taylor
estimates and (1) yield
\[
 |\Delta_h^2f(x)|
 \leq C\frac{h^2}{|x|\log(e/|x|)}
 \leq C\frac{h}{\log(e/h)}=o(h).                              \tag{2}
\]
On compact regions separated from zero, the same second difference is
$O(h^2)$.  The transition regions of $\chi$ are such compact regions.

It remains to treat $|x|<2h$, for which all relevant arguments have
absolute value at most $3h$.  In the core, the second difference of the
linear function $y\mapsto a(h)y$ vanishes.  Therefore its absolute value is
bounded by a fixed multiple of
\[
 hR(h),\qquad
 R(h)=\sup_{0\leq r\leq3h}\frac rh\,|a(r)-a(h)|.             \tag{3}
\]
We claim $R(h)\to0$.  Given $0<\delta<1$, the part $r\leq\delta h$
is at most $2\delta$.  On $\delta h\leq r\leq3h$, write $r=ch$.
Uniformly for $c\in[\delta,3]$,
\[
 \log\log\frac e{ch}-\log\log\frac e h
 =\log\!\left(1-\frac{\log c}{\log(e/h)}\right)
 \longrightarrow0.
\]
Since sine is Lipschitz, the remaining part of (3) tends to zero.  First
choose $\delta$ small and then $h$ small.  Thus (3) is $o(h)$.

Combining this with (2), and taking the monotone supremum over $|h|\leq t$,
proves $\Omega(t)=o(t)$.
\end{proof}

\section{Weak convexity of both sides of the graph}

\begin{lemma}\label{lem:graph}
For every $\varepsilon>0$,
\[
 \gamma_A(\varepsilon),\ \gamma_B(\varepsilon)
 \leq \frac12\Omega(\varepsilon/2).
\]
In particular both exact moduli are $o(\varepsilon)$.
\end{lemma}

\begin{proof}
Take $z_i=(x_i,y_i)\in A$ with $\|z_1-z_2\|\leq\varepsilon$ and set
$x_m=(x_1+x_2)/2$, $y_m=(y_1+y_2)/2$.  Since $y_i\geq f(x_i)$,
\[
 f(x_m)-y_m
 \leq f(x_m)-\frac{f(x_1)+f(x_2)}2
 \leq\frac12\big|\Delta_{(x_1-x_2)/2}^2f(x_m)\big|.
\]
If the midpoint is outside $A$, moving it vertically upward by the positive
quantity on the left reaches $A$.  Its distance to $A$ is therefore at most
$\Omega(\varepsilon/2)/2$.  This proves the estimate for $A$.

For $B$, reverse the inequalities and move a midpoint vertically downward.
The same absolute second difference gives the identical bound.  Lemma
\ref{lem:zyg} now gives both $o(\varepsilon)$ assertions.  In particular
the displayed common upper majorant is less than $\varepsilon/2$ for all
sufficiently small $\varepsilon$, so both sets are weakly convex.
\end{proof}

Closedness and the identity $A=\cl(\operatorname{int}A)$ follow immediately
from the continuity of $f$.

\section{Failure of every nonzero normal}

The Fr\'echet normal cone at $z\in A$ is
\[
 N_A(z)=\left\{p:
 \langle p,w-z\rangle\leq o(\|w-z\|)
 \text{ as }w\to z, w\in A\right\}.
\]
Suppose $p=(u,v)\in N_A(0,0)$.  Testing on the vertical points $(0,s)$,
$s>0$, gives $v\leq0$.

If $v=0$, testing on boundary points $(t,f(t))$ with both signs of $t$
forces $u=0$.  Suppose instead that $v<0$, and put $q=-u/v$.  On boundary
points, after dividing by the positive number $-v$, the normal inequality
implies
\[
 f(t)-qt\geq-o(|t|).                                             \tag{4}
\]
For $t>0$, the quotient
\[
 \frac{f(t)}t=\sin\!\left(\log\log\frac e t\right)
\]
has subsequences converging to both $-1$ and $1$.  Applying (4) along a
subsequence converging to $-1$ gives $q\leq-1$.  For $t<0$, division of
(4) by the negative number $t$ reverses the inequality; along a subsequence
for which $f(t)/t\to1$, it gives $q\geq1$.  This is impossible.
Consequently $p=0$, so $N_A(0,0)=\{0\}$.

Every proximal normal in a Hilbert space is a Fr\'echet normal.  Thus there
is also no nonzero proximal normal at the origin, completing the proof of
Theorem \ref{thm:main}.

\section{Interpretation and verification}

The construction separates the two regularity scales at the heart of the
conjecture.  The estimate $\Delta_h^2f=o(h)$ makes both sides of the graph
asymptotically convex in precisely the required midpoint sense, while the
first-order slopes still oscillate through the whole interval $[-1,1]$.
The hypotheses therefore control symmetric second differences but do not
force a supporting first-order direction.

The accompanying script
\texttt{code/verify\_second\_difference.py} samples the normalized second
differences and the oscillating secants.  It is only a numerical sanity
check; the proof is Lemmas \ref{lem:zyg} and \ref{lem:graph}.

\begin{remark}
The source paper's stronger theorem assumes a quadratic bound on the modulus.
Here the decay is only of order comparable to
$\varepsilon/\log(e/\varepsilon)$ along the difficult scales.  Hence the
example does not contradict that theorem; it shows why the conjectured
replacement by the qualitative condition $o(\varepsilon)$ is too weak.
\end{remark}

\clearpage
\section*{Source context}

The crop below is from page 13 of the locally rendered source
arXiv:1007.0162.  It includes the quadratic theorem, the authors' epilogue,
and Conjecture 6.1.

\begin{center}
\includegraphics[width=0.94\textwidth,trim=0 0 0 86,clip]{figures/open_problem_crop.png}
\end{center}

\begin{thebibliography}{9}
\bibitem{BR}
M.~V. Balashov and D.~Repov\v{s},
\emph{Weakly convex sets and modulus of nonconvexity},
arXiv:1007.0162.
\end{thebibliography}

\end{document}
